% META: {"id":"TSPE-GEOMETRIE-ESPACE-RE-C12","chapitre":"TSPE-GEOMETRIE-ESPACE","type_objet":"remediation","capacites_codes":["C12"],"status":"approved","origin":"generated","status_history":["generated","audited","accepted","approved"],"release_acceptance":"RELEASE_OWNER_BATCH_ACCEPTANCE","acceptance_closure_digest":"sha256:9b3ccf9a81c5520fbb7e03b7b2d3e7bf2057a02834b6d908bf3be5f49f3b553f","acceptance_receipt_digest":"sha256:4fad86f0282e0fa4e0419cca1ad6379ae7af5a280c04a4a90880f90a1c044242"}

%---------------------------------------------------------------
% FICHE DE REMEDIATION — C12 : representation parametrique d'une droite
% Obstacle : tester l'appartenance coordonnee par coordonnee
%---------------------------------------------------------------

\subsection*{Erreur fréquente}

\erreurFrequente{
Pour vérifier qu'un point appartient à une droite, trouver $t$ avec la
première équation et s'arrêter là.
}

\textbf{Pourquoi c'est faux.} Les trois équations partagent le \emph{même}
paramètre. Un $t$ qui convient pour l'abscisse doit convenir aussi pour
l'ordonnée et la cote. S'il change d'une ligne à l'autre, le point n'est pas
sur la droite.

\contreexemple{
Soit $d:\ x=1+2t,\ y=t,\ z=3-t$. Pour $M(5;2;0)$~: la première équation donne
$t=2$, la deuxième aussi. Mais la troisième donnerait $0=3-2=1$~: faux. $M$
n'appartient pas à $d$, alors que deux équations sur trois étaient satisfaites.
}

\subsection*{À retenir}

Une droite admet une infinité de représentations paramétriques --- changer de
point de base ou multiplier le vecteur directeur en donne d'autres. Deux
représentations décrivent la même droite si les vecteurs directeurs sont
colinéaires et si un point de l'une appartient à l'autre.

\subsection*{Méthode corrective}

\begin{enumerate}
  \item Résoudre la première équation en $t$.
  \item Reporter cette valeur dans les \emph{deux} autres.
  \item Conclure seulement si les trois sont vérifiées simultanément.
\end{enumerate}

%---------------------------------------------------------------

\begin{exercice}{TSPE-GEOESPACE-RE-C12-EX1}{2}{12}
Soit $d$ la droite de représentation paramétrique
$x=1+2t$, $y=t$, $z=3-t$, avec $t\in\mathbb{R}$.
\begin{enumerate}
  \item Donner un point de $d$ et un vecteur directeur.
  \item Le point $M(5;2;0)$ appartient-il à $d$ ?
  \item Le point $N(5;2;1)$ appartient-il à $d$ ?
  \item La droite $d'$ définie par $x=3+4s$, $y=1+2s$, $z=2-2s$ est-elle
        confondue avec $d$ ?
\end{enumerate}
\end{exercice}

% BEGIN-VERIFY
% from sympy import Matrix, symbols, solve, Eq
% t, s = symbols('t s')
% A = Matrix([1,0,3]); u = Matrix([2,1,-1])
% # 2. M(5,2,0) : t=2 satisfait x et y, pas z
% M = Matrix([5,2,0])
% assert solve([Eq((A + t*u)[i], M[i]) for i in range(3)], [t]) == []
% assert (A + 2*u)[0] == M[0] and (A + 2*u)[1] == M[1] and (A + 2*u)[2] != M[2]
% # 3. N(5,2,1) : t=2 convient partout
% N = Matrix([5,2,1])
% assert A + 2*u == N
% # 4. d' : point (3,1,2), directeur (4,2,-2) = 2u -> colineaire, et (3,1,2) est sur d (t=1)
% Ap = Matrix([3,1,2]); up = Matrix([4,2,-2])
% assert Matrix.hstack(u, up).rank() == 1
% assert A + 1*u == Ap
% END-VERIFY

\begin{corrige}{TSPE-GEOESPACE-RE-C12-EX1}
\begin{enumerate}
  \item En prenant $t=0$ on obtient le point $A(1;0;3)$, et les coefficients
        de $t$ donnent le vecteur directeur $\vec{u}(2;1;-1)$.
  \item Non. $5=1+2t$ donne $t=2$, et $2=t$ donne aussi $t=2$~; mais la
        troisième équation exigerait $0=3-2=1$, ce qui est faux. Deux
        équations sur trois ne suffisent pas.
  \item Oui. Pour $t=2$~: $x=5$, $y=2$ et $z=3-2=1$. Les trois équations sont
        vérifiées avec la même valeur du paramètre.
  \item Oui. $\vec{u'}(4;2;-2)=2\vec{u}$ donc les directions sont colinéaires,
        et le point $(3;1;2)$ de $d'$ s'obtient sur $d$ pour $t=1$. Les deux
        représentations décrivent la même droite.
\end{enumerate}
\end{corrige}
